\begin{explanationbox}
\textbf{\textcolor{CExplanation}{一句话直觉}}

椭圆面积可以看作单位圆盘面积经过两个方向分别伸缩得到。具体来说，将单位圆盘 $x^2+y^2\le 1$ 沿 $x$ 轴伸缩 $a$ 倍、沿 $y$ 轴伸缩 $b$ 倍，就得到椭圆区域 $\frac{x^2}{a^2}+\frac{y^2}{b^2}\le 1$，面积从 $\pi$ 变为 $\pi ab$。

\medskip
\textbf{\textcolor{CExplanation}{核心拆解}}
\begin{description}[style=nextline, leftmargin=2.8em, labelsep=0.5em, itemsep=0.45em]
    \item[\textbf{要点一（伸缩变换基础）：}] 椭圆区域是单位圆盘经过线性伸缩得到的结果，两个方向的伸缩因子分别为 $a$ 和 $b$。
    \item[\textbf{要点二（面积伸缩因子）：}] 平面图形沿两个互相垂直方向分别伸缩时，面积变换因子等于两个方向伸缩因子的乘积，因此椭圆面积为单位圆盘面积的 $ab$ 倍。
\end{description}

\medskip
\textbf{\textcolor{CExplanation}{几何本质（面积与圆关联）}}

椭圆面积 $\pi ab$ 也可以理解为：该椭圆与半径为 $\sqrt{ab}$ 的圆面积相等。这里的 $\sqrt{ab}$ 是“等面积圆”的半径，不是椭圆本身的半轴。

\medskip
\textbf{\textcolor{CExplanation}{代数意义（对称不可分）}}

面积公式 $S=\pi ab$ 中 $a$ 和 $b$ 的地位完全对称，即使 $a<b$，公式仍然成立，不需要区分谁是长半轴、谁是短半轴。

\medskip
\textbf{\textcolor{CExplanation}{考点价值（直接与反解）}}
\begin{itemize}[leftmargin=*, itemsep=0.5em, topsep=0.4em]
    \item \textbf{考法一（直接求面积）：} 给出椭圆方程，先读出两条半轴，再代入 $S=\pi ab$。
    \item \textbf{考法二（反求半轴）：} 已知面积和一条半轴，反求另一条半轴。
\end{itemize}

\medskip
\textbf{\textcolor{CExplanation}{顿悟点}}

\textbf{在半轴乘积 $ab$ 一定时，椭圆无论“胖瘦”如何变化，面积都相同；若 $ab$ 改变，面积也随之改变。}

\medskip
\textbf{\textcolor{CExplanation}{使用场景}}
\begin{itemize}[leftmargin=*, itemsep=0.5em, topsep=0.4em]
    \item \textbf{场景一（基础面积计算）：} 椭圆以标准形式给出，直接代入公式。
    \item \textbf{场景二（参数反求）：} 已知椭圆的面积和一个半轴，求另一个半轴。
    \item \textbf{场景三（结合焦距）：} 已知短半轴与焦距时，先由 $a^2=b^2+c^2$ 求长半轴，再求面积。
\end{itemize}
\end{explanationbox}